\section{方法}\label{sec:bihc-method}
我们先在第~\ref{sec:bihc-revisit} 节回顾 BEM 及其在双调和坐标推导中的应用，然后在第~\ref{sec:bihc-HighOrderBEM} 节介绍基于高阶 BEM 的高阶笼闭式坐标，最后在第~\ref{sec:bihc-deformation} 节讨论变形控制。

\subsection{边界元方法}\label{sec:bihc-revisit}
\paragraph{BEM 概述}
使用 BEM~\cite{Brebbia1978TheBE} 求解偏微分方程时，通常包含以下步骤：
\begin{enumerate}
    \item 将定义域内部的微分方程转换为边界上的积分方程。
    \item 将边界离散为有限大小的边界单元。
    \item 将边界积分方程转化为代数方程。
    \item 求解代数方程，从而得到原偏微分方程的解。
\end{enumerate}
下面我们说明，应用 BEM 可以在二维线性笼~\cite{Weber2012} 和三维线性笼~\cite{thiery2024biharmonic} 上得到闭式双调和坐标。

\paragraph{双调和 Dirichlet 问题}
双调和坐标由有界域 $\Omega$ 上双调和 Dirichlet 问题的解导出（参见~\cite{Weber2012} 式 (4)）：
\begin{equation}\label{equ:bihc-bihPDE}
\begin{aligned}
    &\Delta^2 f(\eta) = 0, \quad \eta \in \Omega, \\
    &f(\xi) = g_1(\xi), \quad \xi \in \partial \Omega, \\
    &\frac{\partial f}{\partial \mn_{\xi}}(\xi) = g_2(\xi), \quad \xi \in \partial \Omega,
\end{aligned}
\end{equation}
其中，$f$ 是双调和函数，\(g_1\) 与 \(g_2\) 分别为给定的 Dirichlet 和 Neumann 边界条件，\(\mn_{\xi}\) 是笼边界 \(\partial\Omega\) 在 \(\xi\) 处的法向量。

\paragraph{边界积分方程}
利用 Green 定理和基本解，可将偏微分方程~\eqref{equ:bihc-bihPDE} 转化为如下边界积分方程：
\begin{equation}\label{equ:bihc-BIE1}
\begin{aligned}
f(\eta) &=
\int\limits_{\xi \in \partial \Omega}
f(\xi) \frac{\partial G_1}{\partial \mn_{\xi}}(\xi, \eta)d\xi -
\int\limits_{\xi \in \partial \Omega}
G_1(\xi, \eta) \frac{\partial f(\xi)}{\partial \mn_{\xi}}d\xi \\
+ &\int\limits_{\xi \in \partial \Omega}
\Delta f(\xi) \frac{\partial G_2}{\partial \mn_{\xi}}(\xi, \eta)d\xi -
\int\limits_{\xi \in \partial \Omega}
G_2(\xi, \eta) \frac{\partial \Delta f(\xi)}{\partial \mn_{\xi}}d\xi,
\end{aligned}
\end{equation}

\begin{equation}\label{equ:bihc-BIE2}
\begin{aligned}
\Delta f(\eta) =
\int\limits_{\xi \in \partial \Omega}
\Delta f(\xi) \frac{\partial G_1}{\partial \mn_{\xi}}(\xi, \eta)d\xi -
\int\limits_{\xi \in \partial \Omega}
G_1(\xi, \eta) \frac{\partial \Delta f(\xi)}{\partial \mn_{\xi}}d\xi,
\end{aligned}
\end{equation}
其中，$G_1$ 和 $G_2$ 分别是二维调和方程与双调和方程的基本解。具体写为：
$G_1(\xi,\eta)=-\frac{1}{2 \pi}\ln\|\xi-\eta\|$，
$G_2(\xi,\eta)=-\frac{\|\xi-\eta\|^2}{8 \pi}(\ln\|\xi-\eta\|-1)$。
%
在积分方程中，由于~\eqref{equ:bihc-bihPDE} 中已给定 \(f\) 与 \(\partial f/\partial \mn\)，未知函数仅包含 \(\Delta f\) 和 \(\partial \Delta f/\partial \mn\)。
当我们在边界上求得 \(\Delta f\) 与 \(\partial \Delta f/\partial \mn\) 的全部未知函数值后，将其代入~\eqref{equ:bihc-BIE1}，即可得到域内任意点的 $f$ 值，即~\eqref{equ:bihc-bihPDE} 的解。
%
为在边界上求解 \(\Delta f\) 与 \(\partial \Delta f/\partial \mn\)，需要在~\eqref{equ:bihc-BIE1} 与~\eqref{equ:bihc-BIE2} 中将 \(\eta\) 移到边界上（记作 \(\eta^{\text{b}}\)）。

\paragraph{线性笼上的双调和坐标}
一般情况下，我们将积分方程离散化为代数方程，并进一步数值求解。
%
离散过程包括两步。
第一步是将边界划分为有限大小的边界单元。对于二维多边形笼~\cite{Weber2012} 与三维三角网格笼~\cite{thiery2024biharmonic}，边界是分段线性的，因此可自然离散为分段线性单元。
%
第二步是在每个边界单元内用插值逼近边界函数。
在每个单元中，\(f\) 与 \(\partial f/\partial \mn\) 分别用分段线性函数与分段常函数表示。未知量 \(\Delta f\) 与 \(\partial \Delta f/\partial \mn\) 也分别用分段线性函数与分段常函数近似。
%
将这些函数表示为边界单元节点处函数值后，即可通过在不同边界点上评估 \(\eta'\) 建立代数方程组。
%
在代数方程组求解方面，\citet{Weber2012} 将代数约束~\eqref{equ:bihc-BIE2} 代入~\eqref{equ:bihc-BIE1}，形成线性系统并求解。
在三维情形中，\citet{thiery2024biharmonic} 同时最小化两类约束，从而得到更平稳的变形结果。

\begin{figure}[t]
  \centering
  \begin{overpic}[width=0.99\linewidth]{diagram}
    {
 \put(17,12){\small $\eta$}
  \put(29,29){\small $\mc_i$}
   \put(87,27){\small $\mg_1$}
    }
  \end{overpic}
  \vspace{-3mm}
  \caption{
  高阶边界元方法示意。曲边边界由一系列三次单元 $\mc_i$ 构成，$\mg_1$ 在其上为三次多项式函数。
  }
  \label{fig:bihc-schematic}
\end{figure}

\subsection{多项式双调和坐标}\label{sec:bihc-HighOrderBEM}

为实现曲边到曲边的变形，我们采用高阶边界元方法~\cite{Fischer1999ApplicationOH,Frijns2000ImprovingTA,gravenkamp2019high}，这需要高阶单元和高阶形函数。


\subsubsection{重写边界积分方程}\label{sec:bihc-reformulating-bie}
\paragraph{高阶单元}
输入二维笼由 $N_c$ 条多项式曲线组成。
这些曲线在边界积分方程中彼此独立。
%
第 $i$ 个边界单元是一条 $m$ 次多项式曲线，参数化为 \(\mc_i(t): [0, 1] \to \mathbb{R}^2\)。
该曲线 \(\mc_i(t)\) 可表示为 $m$ 次 \Bezier 曲线：
\begin{equation}
    \mc_i(t)=\sum_{j=0}^m \mc_{j}^i B_j^m(t),
\end{equation}
其中 \(B_j^m\) 表示 Bernstein 基函数，\(\mc_{j}^i\) 表示对应控制点。

\paragraph{高阶边界条件}
通常我们假设 \(\mg_1\) 在单元 \(\mc_i(t)\) 上是 $n$ 次多项式函数：
\begin{equation}\label{equ:bihc-g}
\begin{aligned}
    \mg_1(t)=\sum_{j=0}^n \mg_{1,j}^i B_j^n(t), \,\, \mg_2(t)=s_i\overbrace{\frac{\|\mg_1'(t)\|}{\|\mc_i'(t)\|}}^{\sigma} \overbrace{\frac{\mg_1'(t)^{\perp}}{\|\mg_1'(t)^{\perp}\|}}^{\bar{\mn}}=s_i\frac{\mg_1'(t)^{\perp}}{\|\mc_i'(t)\|},
    \end{aligned}
\end{equation}
其中 \(\mg_{1,j}^i\) 为控制点，且 \((a, b)^\perp = (b, -a)\)。
这里我们将缩放因子 \(\sigma\) 乘以单位法向来施加 Neumann 边界条件。若简单将 \(\sigma\) 设为 1，会出现不期望项 \(\|\mg_1'(t)^\perp\|\)。因此我们采用与 Green 坐标相同的 \(\sigma\) 形式，从而保证与~\cite{MichelThiery2023} 相近的边界行为。
此外我们引入缩放因子 \(s_i\)，其默认取值为 1。更复杂的关于 \(s_i\) 的 Neumann 边界条件配置见第~\ref{sec:bihc-deformation} 节。
%
由于 \(\mg_1(t)\) 意味着变形后笼边是 $n$ 次多项式曲线，我们取 \(n \geq m\) 以保证变形连续性。

\paragraph{高阶形函数}
我们使用高阶形函数来近似边界上未知函数 \(\Delta f\) 和 \(\partial \Delta f/\partial \mn\)。
通常，在 \(\mc_i(t)\) 上用 $k$ 次多项式近似 \(\Delta f\)，用 $(k-1)$ 次多项式近似 \(\partial \Delta f/\partial \mn\)：
\begin{equation}\label{equ:bihc-Df}
\begin{aligned}
    \Delta f(t)=\sum_{j=0}^k \mh_{1,j}^i B_j^k(t), \frac{\partial \Delta f}{\partial \mn}(t)=\sum_{j=0}^{k-1} \mh_{2,j}^i B_j^{k-1}(t),
    \end{aligned}
\end{equation}
其中 \(\mh_{1,j}^i\) 与 \(\mh_{2,j}^i\) 为控制点。
我们取 \(k \geq n\) 以保证可重现仿射变换。

\paragraph{方程重写}
注意 $\frac{\partial G}{\partial \mn}=\nabla_1 G(\mc(t),\eta)\cdot \mn_{c(t)}$，且 $\mn_{\mc(t)} d\xi=\mn_{\mc(t)}\|\mc'(t)\|dt=\mc'(t)^{\perp}dt$。
将 \eqref{equ:bihc-g} 与 \eqref{equ:bihc-Df} 代入 \eqref{equ:bihc-BIE1} 和 \eqref{equ:bihc-BIE2}，得到离散化后的边界积分方程：
\begin{equation}\label{equ:bihc-dis_Int1}
\begin{aligned}
    f(\eta) = \sum_{i=1}^{N_c} &\left(\sum_{j =0}^n \phi_{i,j}(\eta) \mg_{1,j}^i +  \sum_{j =0}^{n-1}\psi_{i,j}(\eta) s_i (\mg_{2,j}^{i})^{\perp} \right.\\
&+\left. \sum_{j =0}^k \tilde{\phi}_{i,j}(\eta) \mh_{1,j}^i +  \sum_{j =0}^{k-1}\tilde{\psi}_{i,j}(\eta) (\mh_{2,j}^{i})^{\perp}\right),
\end{aligned}
\end{equation}
\begin{equation}\label{equ:bihc-dis_Int2}
   \Delta f(\eta) = \sum_{i=1}^{N_c} \left(\sum_{j =0}^k  \phi_{i,j}(\eta) \mh_{1,j}^i + \sum_{j =0}^{k-1} \psi_{i,j}(\eta) (\mh_{2,j}^{i})^{\perp}\right),
\end{equation}
其中 $\mg_{2,j}^i=n(\mg_{1,j+1}^i-\mg_{1,j}^i)$，且
\begin{equation}\label{equ:bihc-basis_inte}
\begin{aligned}
   & \phi_{i,j}(\eta)=\int\limits_0^1\nabla_1 G_1(\mc_i(t),\eta)\cdot \mc_i'(t)^{\perp}B^n_j(t)\,dt,\\
   &\psi_{i,j}(\eta)=\int\limits_0^1 G_1(\mc_i(t),\eta) B^n_j(t)\,dt,\\
   &\tilde{\phi}_{i,j}(\eta)=\int\limits_0^1\nabla_1 G_2(\mc_i(t),\eta)\cdot \mc_i'(t)^{\perp}B^n_j(t)\,dt,\\
   &\tilde{\psi}_{i,j}(\eta)=\int\limits_0^1 G_2(\mc_i(t),\eta) \|\mc_i'(t)\|^2 B^n_j(t)\,dt,
 \end{aligned}
\end{equation}
其中 $\nabla_1 G_1(\xi, \eta) = \frac{\xi - \eta}{2\pi \|\xi - \eta\|^2}$，$\nabla_1 G_2(\xi, \eta) = \frac{(\xi - \eta)(2 \log(\|\xi-\eta\|)-1)}{8\pi}$。
%
\paragraph{积分及其导数的解析计算}
与 \(\phi_{i,j}\) 和 \(\psi_{i,j}\) 相关的项，正是多项式 Green 坐标中的项，已由~\cite{MichelThiery2023} 针对多边形笼、由~\cite{liu2024polygc} 针对高阶笼给出。
这里我们先给出 \(\tilde{\phi}_{i,j}\) 与 \(\tilde{\psi}_{i,j}\) 的计算方法，再推导它们的导数。

对于 \(\tilde{\phi}_{i,j}\)，注意
\(\nabla_1 G_2(\mc_i(t),\eta) \cdot \mc_i'(t)^{\perp} = (-\eta \cdot \mc_i'(t)^{\perp})(2 \log(\| c(t) - \eta\|) - 1)\)。
\(\tilde{\phi}_{i,j}\) 中的积分项可写为“多项式函数与对数项”的乘积。为简洁起见，我们可先关注以下项：
\begin{equation}\label{equ:bihc-cal_termsingle}
\begin{aligned}
    &\int_0^1 t^j \log(\|\mc_i(t) - \eta \|)\, dt \\
    &= \frac{\log(\|\mc^i_{m} - \eta\|)}{j+1} + \frac{1}{j+1}\int_0^1 \frac{t^{j+1}(c(t)-\eta)\cdot c'(t)}{\|c(t)-\eta\|^2}\, dt,
\end{aligned}
\end{equation}
其中 \(j = 0, \dots, n\)。问题被转化为计算
$F_{h,2}^{c,\eta}=\int_0^1\frac{t^h}{\|c(t)-\eta \|^2}\,dt,\ h=j+1,\dots,j+2m$。
类似地，\(\tilde{\psi}_{i,j}\) 中积分项也可写为多项式与对数项乘积，其计算方式与 \(\tilde{\phi}_{i,j}\) 相同，但多项式次数更高，此时 \(j = 0, \dots, 4m + n - 2\)。

导数与积分具有相似形式。对 \(\tilde{\phi}_{i,j}'\) 来说，非平凡项仅涉及 \eqref{equ:bihc-cal_termsingle} 的导数。
实际上，对 \eqref{equ:bihc-basis_inte} 中四类积分求导时，非平凡项仅包含：
$
F_{h,4}^{c,\eta} := \int_0^1 \frac{t^h}{\|c(t) - \eta \|^4} \, dt.
$

因此，只需计算 \(F_{h,2}^{c,\eta}\) 与 \(F_{h,4}^{c,\eta}\)。\citet{liu2024polygc} 给出了如下引理：
\begin{lemma}
对于 \(p(t) = \prod_{i=1}^{n} (t - \omega_i)^{n_i}\)，有
\[
\int_0^1 \frac{t^h}{p(t)} dt = \sum_{i=1}^{n} \text{Res}\left(\frac{w^m \left( \log\left(1 - \frac{1}{w}\right) + \sum_{k=1}^h \frac{1}{k w^k} \right)}{p(w)}, \omega_i\right).
\]
\end{lemma}
我们将该引理用于计算 \(F_{h,2}^{c,\eta}\) 与 \(F_{h,4}^{c,\eta}\)，即令
$p(t)=\|c(t) - \eta \|^2$ 与 $p(t)=\|c(t) - \eta \|^4$。
这涉及求解方程 \(\|c(t)-\eta\|^2=0\)。当 \(c(t)\) 的次数小于 5 时，该方程有解析解。

\subsubsection{边界积分方程求解}
\paragraph{奇异积分}
为得到代数方程组，我们在边界上采样 \(\eta^{\text{b}}\)。
%
此时，上述积分求值涉及奇异积分计算。关于该问题，在边界元方法中已有大量研究~\cite{Fischer1999ApplicationOH}。我们遵循~\cite{Weber2012} 的建议：当 \(\eta\) 趋近边界时，这些积分的极限存在。
\tr{在本问题中，上述解析方法无法直接计算该极限，因为每一项 \(F_{h,2}^{c,\eta}\) 都趋于无穷大。我们在附录~\ref{sec:bihc-appendix} 中单独给出该极限的显式表达。另一种实用策略是将点 \(\eta^{\text{b}}\) 略微向笼内部偏移，再直接计算积分。实验表明，这一做法同样可得到准确结果。}

\paragraph{矩阵约束}
将边界上的 \(\eta^{\text{b}}\) 代入 \eqref{equ:bihc-dis_Int1} 与 \eqref{equ:bihc-dis_Int2}，可得如下矩阵约束：

\begin{equation}\label{equ:bihc-matrix1}
    M_nG_1=\Phi_nG_1+\Psi_n G_2+\bar{\Phi}_k H_1+\bar{\Psi}_k H_2,
\end{equation}
\begin{equation}\label{equ:bihc-matrix2}
    M_k H_1=\Phi_k H_1+\Psi_k H_2.
\end{equation}
其中，\(G_1\)（或 \(G_2\)）是由 \(\mg_{1,j}^i\)（或 \(s_i (\mg_{2,j}^{i})^{\perp}\)）组成的 \(N_c\cdot n\) 维列向量；
\(H_1\)（或 \(H_2\)）是由 \(\mh_{1,j}^i\)（或 \((\mh_{2,j}^{i})^{\perp}\)）组成的 \(N_c\cdot k\) 维列向量。
注意，由于连续性，相邻边共享公共的 \(\mg_{1,0}^i\) 与 \(\mh_{1,0}^i\)（或 \(\mg_{1,n}^i\) 与 \(\mh_{1,k}^i\)）。
式~\eqref{equ:bihc-matrix1} 和 \eqref{equ:bihc-matrix2} 中矩阵的第 \(l\) 行，对应于在第 \(l\) 个采样点处对 \eqref{equ:bihc-dis_Int1} 与 \eqref{equ:bihc-dis_Int2} 的计算。
\(M_n\) 与 \(M_k\) 收集了位于对应边界单元上的 \(\eta^{\text{b}}\) 处，$n$ 次与 $k$ 次 Bernstein 基函数系数。
\(\Phi_n\) 按照 \(G_1\) 的顺序收集 \eqref{equ:bihc-dis_Int1} 中的 \(\phi_{i,j}(\eta^{\text{b}})\)。
\(\Psi_n,\bar{\Phi}_k,\bar{\Psi}_k,\Phi_k\) 和 \(\Psi_k\) 定义方式相同。


\paragraph{近似代数方程}
求解代数方程组有两种方式：一种严格施加拉普拉斯约束，另一种较宽松地施加该约束。
为简化说明，我们只讨论 \(k=n\) 的情形；\(k>n\) 的情况同理。
此时，方程组中的 8 个矩阵可简化为 5 个。由此可得到我们正则化双调和坐标的表达：
\begin{equation}\label{equ:bihc-final_BHC}
    \begin{aligned}
f(\eta) &= \big[ \phi(\eta) + (\bar{\phi}(\eta) C_L + \bar{\psi}(\eta) C_D)(M_n - \Phi_n) \big] G_1 \\
&\quad + \big[ \psi(\eta) - (\bar{\phi}(\eta)) C_L + \bar{\psi}(\eta) C_D)\Psi_n \big] G_2  \\
&:=\alpha(\eta)G_1+\beta(\eta)G_2
\end{aligned}
\end{equation}
其中，\(\phi(\eta)\) 表示在 \(\eta\) 处由 \(\phi_{i,j}\) 取值组成的行向量；\(\bar{\phi}(\eta)\)、\(\bar{\psi}(\eta)\) 和 \(\psi(\eta)\) 同理。
\(\alpha(\eta)\) 与 \(\beta(\eta)\) 共同构成我们的双调和坐标。
在~\cite{Weber2012} 提出的第一种方法（记作 $\text{BiHC}_1$）中，仅优化 \eqref{equ:bihc-dis_Int1}，同时严格满足 \eqref{equ:bihc-dis_Int2}。
相应矩阵为：\(C_L = (\bar{\Phi}_n + \bar{\Psi}_n A)^{-1}\)、\(C_D = A (\bar{\Phi}_n + \bar{\Psi}_n A)^{-1}\)，其中 \(A = \Psi_n^{-1} (M_n - \Phi_n)\)。
%
在~\cite{thiery2024biharmonic} 的第二种方法（记作 $\text{BiHC}_{1,2}$）中，同时优化 \eqref{equ:bihc-dis_Int1} 与 \eqref{equ:bihc-dis_Int2}。
对应矩阵通过下式确定：
\[
\begin{pmatrix}
C_L \\
C_D
\end{pmatrix} = \left( E^T E + F^T F \right)^{-1} E^T,
\]
其中 \(E = (\bar{\Phi}_n; \bar{\Psi}_n)\)，\(F = (\Phi_n - M_n; \Psi_n)\)。
实验表明，当变形幅度较小时，$\text{BiHC}_{1}$ 与 $\text{BiHC}_{1,2}$ 结果接近；而在更具挑战性的变形下，$\text{BiHC}_{1,2}$ 的变形结果优于 $\text{BiHC}_{1}$（见图~\ref{fig:bihc-linear-solving}）。这一观察与~\cite{thiery2024biharmonic} 的结论一致：联合优化拉普拉斯与 Dirichlet 约束可提供更多自由度。


\subsection{变形控制}\label{sec:bihc-deformation}

\paragraph{提高精度}
为提高数值解精度，可采用三种方法：
\begin{enumerate}
    \item 增加边界单元数量，以扩展离散近似空间。
    \item 增大 \(k\) 的次数。与方法 (1) 类似，这会得到更大的近似空间和更多自由度。
    \item 增加采样点 \(\eta^\text{b}\) 数量，形成超定方程组，使积分方程在最小二乘意义下于更多点成立，从而使变形结果更光滑。
\end{enumerate}
实验表明，方法 (1) 与方法 (2) 结果相近。
在实践中，我们将每条笼边细分为 4 个单元，并在每个单元参数域上均匀采样 \(2n\) 个点构建方程组。

\begin{figure}[t]
  \centering
  \begin{overpic}[width=1.0\linewidth]{diff_solver}
    {
    \put(7,-3.5){\small 初始姿态}
    \put(40,-3.5){\small $\text{BiHC}_{1}$}
 \put(78,-3.5){\small $\text{BiHC}_{1,2}$}
    }
  \end{overpic}
  \caption{
  两种代数方程求解方法的比较。当仅对裤子右侧进行变形时，$\text{BiHC}_{1}$ 在左腿区域引入了扰动，而 $\text{BiHC}_{1,2}$ 保持了无扰动变形。
  }
  \label{fig:bihc-linear-solving}
\end{figure}


\paragraph{变形能量}
在~\eqref{equ:bihc-g} 中，不同的 \(s_j\) 选择对应不同的 Neumann 边界条件\tr{ \(g_2\)。其中 \(g_2\) 表示双调和函数在边界上的导数，刻画了映射在边界附近的拉伸行为。}
\tr{
令 \(s_j=1\) 时，得到与~\cite{MichelThiery2023} 相同的边界条件。不同于其公式，这类边界条件可在双调和函数空间中被精确满足。我们观察到，该选择在边界附近呈现与 Green 坐标相似的保角行为，并在多组样例中取得了良好结果。}
然而，对于偏离保角较大的笼变形，这种边界附近的局部保角性可能导致笼内部产生更大的剪切。
在这种情况下，需要考虑笼的全局变形。
按照~\cite{Weber2012} 与~\cite{thiery2024biharmonic} 的做法，我们给出两种根据边界位置自动推导边界导数的方法。
由于 \(s_i<0\) 在变形时可能导致边界附近退化与折叠，我们约束 \(s_i>0\)。
我们在 \(\partial\Omega\) 上采样点约束 $C=\{q_k\}$，以离散两种不同的变形能量：
\begin{enumerate}
\item 尽可能调和（As-Harmonic-As-Possible, AHAP）能量：$\sum_k \|\Delta f(q_k)\|^2$。该能量通过~\eqref{equ:bihc-dis_Int2} 计算。
\item 尽可能仿射（As-Affine-As-Possible, AAAP）能量：$\sum_k \|f'(q_k)\|^2$。该能量通过~\eqref{equ:bihc-final_BHC} 计算。其评估需要积分 \(\phi(\eta),\psi(\eta),\bar{\phi}(\eta),\bar{\phi}(\eta)\) 的导数，而这部分计算已在第~\ref{sec:bihc-reformulating-bie} 节给出。
\end{enumerate}
我们发现，单位法向导数，以及通过 AHAP 与 AAAP 能量得到的法向导数，能够逐步增强内部保角性。三种方法都可产生可用变形，具体取决于应用场景。


\begin{figure}[t]
  \centering
  \begin{overpic}[width=1.0\linewidth]{scale}
    {
    \put(6,-3.5){\small 初始姿态}
    \put(35,-3.5){\small 单位}
    \put(58,-3.5){\small AHAP}
 \put(85,-3.5){\small AAAP}
    }
  \end{overpic}
  \caption{
  不同缩放因子的变形比较，包括单位缩放和通过最小化变形能量优化得到的缩放。
  }
  \label{fig:bihc-scaling}
\end{figure}

\begin{figure}[t]
  \centering
  \begin{overpic}[width=0.99\linewidth]{diff_order_pangxie}
    {
    \put(2,-3.5){\small \textbf{(a)} 初始姿态}
     \put(26,-3.5){\small \textbf{(b)} $\omega=0$}
      \put(55,-3.5){\small \textbf{(c)} $\omega=1$}
 \put(80,-3.5){\small \textbf{(d)} $\omega=0.5$}
    }
  \end{overpic}
  \caption{
  加权坐标结果 (d)，其表示为 Green 坐标 (b) 与双调和坐标 (c) 的线性组合。
  }
  \label{fig:bihc-weights}
\end{figure}

\paragraph{不同权重}
在笼变形本身非保角时，仅用双调和坐标无法得到严格保角变形，因为“保角性”和“插值性”本质上相互冲突。
不过，双调和坐标可以让用户在这两者之间进行平衡并按需侧重。
%
在~\eqref{equ:bihc-final_BHC} 中，可将双调和坐标拆分为两部分：\(\{\alpha_c,\beta_c\}\) 与 \(\{\alpha_{\text{i}},\beta_{\text{i}}\}\)，定义如下：
\begin{equation*}
    \begin{split}
    &\alpha_c(\eta)=\phi(\eta),\\
    &\beta_c(\eta)=\psi(\eta),\\
     &\alpha_{\text{i}}(\eta)=(\bar{\phi}(\eta) C_L + \bar{\psi}(\eta) C_D)(M_n - \Phi_n),\\
     &\beta_{\text{i}}(\eta)=- (\bar{\phi}(\eta)) C_L + \bar{\psi}(\eta) C_D)\Psi_n.
    \end{split}
\end{equation*}
其中 \(\{\alpha_c,\beta_c\}\) 对应 Green 坐标并提供保角变形，\(\{\alpha_{\text{i}},\beta_{\text{i}}\}\) 是一个双调和函数项，用于把保角变形后的边界拉回到笼边界。
%
我们定义新坐标 \(\{\alpha_c + w \alpha_{\text{i}},\beta_c+w\beta_{\text{i}}\}\)，用户可调节 \(w\) 以控制变形。
当 \(w=0\) 时，变形是保角的；当 \(w=1\) 时，结果完全插值边界。
通常取中间值 \(w\) 可得到更直观且更满意的变形效果（见图~\ref{fig:bihc-weights}）。


\begin{algorithm}[t]
  \caption{PolyBiharmonicCoordinates：计算多项式双调和坐标（\(\text{BiHC}_1\)）}\label{alg:bihc-encoding}
  \SetKwInput{KwInput}{输入}
  \SetKwInput{KwOutput}{输出}
  \SetKwFunction{FMain}{Main}
  \SetKwFunction{FD}{D}
  \KwInput{\(c\)：由 $N_c$ 条 $n$ 次多项式曲线构成的笼\\
          \quad\quad\quad  \(\eta\)：笼内部一点}
  \KwOutput{$\alpha(\eta)$ 与 $\beta(\eta)$：点 $\eta$ 处的双调和坐标}
  % Algorithm
  \tb{//构建矩阵\;}
  \tr{
  $\{\eta^b\}\gets$ 每条曲线上采样 $n$ 个点\;
  \For{每个 $\eta^b$}{
  计算积分 $\phi_{i,j}(\eta^b),\psi_{i,j}(\eta^b),\tilde{\phi}_{i,j}(\eta^b),\tilde{\psi}_{i,j}(\eta^b)$\;
  组装系统矩阵：$M_n,\Phi_n,\Psi_n,\bar{\Phi}_n,\bar{\Psi}_n$\;
  }
  $A=\Phi_n^{-1}(M_n-\Phi_n),C_L=(\bar{\Phi}_n+\bar{\Psi}_nA)^{-1},C_D=AC_L$\;
  }
  \tb{//计算坐标\;}
    \tr{
  计算积分：$\phi_{i,j}(\eta),\psi_{i,j}(\eta),\tilde{\phi}_{i,j}(\eta),\tilde{\psi}_{i,j}(\eta)$\;
  组装列向量：$\phi(\eta),\psi(\eta),\bar{\phi}(\eta),\bar{\psi}(\eta)$\;
  $\alpha(\eta)=\phi(\eta) + (\bar{\phi}(\eta) C_L + \bar{\psi}(\eta) C_D)(M_n - \Phi_n)$\;
  $\beta(\eta)=\psi(\eta) - (\bar{\phi}(\eta)) C_L + \bar{\psi}(\eta) C_D)\Psi_n$\;
  }
\end{algorithm}

\begin{figure*}[t]
  \centering
  \begin{overpic}[width=0.99\linewidth]{gallery}
    {
    }
  \end{overpic}
  \vspace{-3mm}
  \caption{
  将我们的方法应用于更多多项式曲线笼示例。
  }
  \label{fig:bihc-more-examples}
\end{figure*}
